\documentclass[a4paper,12pt,oneside]{article}

\usepackage[T2A]{fontenc}
\usepackage[utf8]{inputenc}
\usepackage[russian]{babel}

\usepackage{tempora}
\usepackage{microtype}

\usepackage{amsmath,amssymb,amsthm,mathtools}
\usepackage{geometry}
\geometry{top=25mm,left=25mm,right=20mm,bottom=25mm}

\usepackage{enumitem}
\setlist[enumerate]{leftmargin=*, itemsep=0.75\baselineskip, topsep=0.35\baselineskip}
\setlist[itemize]{leftmargin=*, itemsep=0.15\baselineskip, topsep=0.2\baselineskip}

\setlength{\parindent}{1.25em}
\setlength{\parskip}{0pt}

\newcommand{\bb}{\mathfrak{B}}
\newcommand{\NOT}[1]{\overline{#1}}
\newcommand{\BigDisp}[1]{%
\[
{\Large\displaystyle #1}
\]
}
\newcommand{\DefItem}[2]{\item \textbf{#1}\\#2}

\begin{document}

\begin{center}
    {\Large\bfseries Теоретический минимум}\\[2mm]
    {\normalsize 2 модуль \hfill}
\end{center}
\vspace{2mm}
\hrule
\vspace{4mm}

\begin{enumerate}[label=\arabic*.]

\DefItem{Определение высказывания}{Утверждение, о котором можно сказать, что оно истинно или ложно.}

\DefItem{Арность операции (функции)}{Число аргументов данной операции (функции).}

\DefItem{Булева алгебра (определение + аксиомы/законы только перечислить)}{Множество $B$, в которое входят элементы $0$ и $1$, на котором заданы бинарные операции конъюнкции и дизъюнкции, и унарная операция отрицания, и выполняются следующие аксиомы для всех $a,b,c$ из $B$:
\begin{itemize}
\item закон коммутативности;
\item закон ассоциативности;
\item закон дистрибутивности;
\item закон тождества;
\item закон дополнения.
\end{itemize}}

\DefItem{Булево множество}{Множество $\bb$, состоящее только из элементов $0$ и $1$.}

\DefItem{Закон единственности дополнения}{Дополнение произвольного элемента $x$ единственным образом определяется его свойствами:
\BigDisp{x+\NOT{x}=1;\qquad x\NOT{x}=0.}}

\DefItem{Закон инволюции}{Для любого $a$ из булевой алгебры верно:
\BigDisp{\NOT{\NOT{a}}=a.}}

\DefItem{Теорема склеивания}{Для $x$ и $y$ из булевой алгебры верно:
\[
\begin{aligned}
xy + x\overline{y} &= x,\\
(x+y)(x+\overline{y}) &= x.
\end{aligned}
\]}

\DefItem{В чем отличие функции и формулы?}{У функции может быть множество формул.}

\DefItem{Число булевых функций $n$ аргументов}{\BigDisp{2^{2^{n}}.}}

\DefItem{ДНФ}{Дизъюнкция выражений, которые либо:
\begin{itemize}
\item отдельный аргумент (возможно с инверсией);
\item простая конъюнкция некоторых аргументов.
\end{itemize}}

\DefItem{КНФ}{Конъюнкция выражений, которые либо:
\begin{itemize}
\item отдельный аргумент (возможно с инверсией);
\item простая дизъюнкция некоторых аргументов.
\end{itemize}}

\DefItem{Минтерм}{Минтермом $n$ переменных называется такая конъюнкция их, в которую каждая переменная входит один раз в прямой или инверсной форме.\\
\textbf{Утв:} Минтерм — это булева функция, которая принимает единичное значение только на одном наборе значений переменных.}

\DefItem{СДНФ (совершенная дизъюнктивная нормальная форма)}{Представление функции в виде дизъюнкции минтермов $n$ аргументов.}

\DefItem{Макстерм}{Макстермом $n$ переменных называется такая дизъюнкция их, в которую каждая переменная входит один раз в прямой или инверсной форме.\\
\textbf{Утв:} Макстерм — это булева функция, принимающая единичные значения на всех наборах аргументов, кроме одного.}

\DefItem{СКНФ}{КНФ, которая не имеет одинаковых дизъюнкций и все они полные.}

\DefItem{Теорема разложения для ДНФ/КНФ}{\textbf{ДНФ:} Всякую булеву функцию можно представить в виде:
\BigDisp{f=(\NOT{x}\,f|_{x=0})+(x\,f|_{x=1}).}
\textbf{КНФ:} Всякую булеву функцию можно представить в виде:
\BigDisp{f=(x+f|_{x=0})(\NOT{x}+f|_{x=1}).}}

\DefItem{Импликанта функции $f$}{Функция, все минтермы которой входят в множество минтермов функции $f$.}

\DefItem{Простая импликанта}{Импликанта, состоящая из одной простой конъюнкции, такой, что никакая её часть не является импликантой функции.\\
\textbf{Вариант 2:} удаление из неё \textbf{любого} литерала приводит к тому, что она перестает быть импликантой функции.}

\DefItem{Число импликант функции}{\BigDisp{2^{k},\ \text{где } k\ \text{— число минтермов}.}}

\DefItem{Минимизация булевой формулы}{Нахождение наименьшего числа простых импликант, которые в дизъюнкции описывают исходную функцию.\\
\textbf{Методы минимизации:} алгебраический, метод Петрика, метод Квайна—Мак-Класки, карты Карно.}

\DefItem{Сокращенная ДНФ}{Дизъюнкция всех простых импликант.}

\DefItem{Тупиковая ДНФ}{Дизъюнкция простых импликант, ни одна из которых не является лишней (импликанта, при удалении которой из формулы функция не меняется).}

\DefItem{Минимальная ДНФ}{Тупиковая ДНФ, которая содержит минимальное число вхождений переменных.}

\DefItem{Полином Жегалкина. Методы построения}{Полином с коэффициентами $0$ и $1$, где произведение представлено операцией конъюнкции, а сложение — сложением по модулю $2$.}

\DefItem{Суперпозиция функций}{Функция, полученная из некоторого множества функций путём подстановок одной функции в другую и/или отождествления переменных.}

\DefItem{Подстановка функции $g$ в функцию $f$}{Замена $i$-го аргумента функции $f$ значением функции $g$.}

\DefItem{Отождествление переменных в функции $f$}{Постановка $i$-го аргумента функции вместо $j$-го.}

\DefItem{Замкнутое множество функций}{Такое множество, что любая функция алгебры логики, выражаемая с помощью содержащихся в множестве функций, уже содержится в этом множестве.\\
\textbf{Вариант 2:} множество $A$ называется замкнутым классом, если замыкание $A$ совпадает с самим $A$: $[A]=A$.}

\DefItem{Замыкание множества функций $A$}{Множество всех суперпозиций функций из множества $A$.\\
\textbf{Вариант 2:} замыканием множества функций называется минимальное замкнутое подмножество всех функций, содержащее данное множество функций.\\
\textbf{Вариант 3:} замыканием $A$ называется множество всех функций алгебры логики, которые можно выразить формулами над $A$.}

\DefItem{Полная система функций}{Множество функций, для которого замыкание совпадает с множеством всех булевых функций.}

\DefItem{Безызбыточная полная система функций / Избыточная полная система функций}{Полная система функций, которая перестает быть таковой после исключения из неё любой функции.}

\DefItem{Самодвойственная функция. Число самодвойственных функций $n$ аргументов}{Функция, которая на противоположных наборах возвращает противоположные значения. Число самодвойственных функций:
\BigDisp{2^{2^{\,n-1}}.}}

\DefItem{Линейная функция. Число линейных функций $n$ аргументов}{Функция, которая не содержит конъюнкций в своем представлении в виде полинома Жегалкина. Число линейных функций:
\BigDisp{2^{n+1}.}}

\DefItem{Монотонная функция}{Функция, которая на сравнимых наборах не убывает.}

\DefItem{Сравнимые наборы}{Наборы $a$ и $b$ сравнимы, если для любого $i$ выполнено $a_i \le b_i$.}

\DefItem{Функции, сохраняющая константу. Число таких функций $n$ аргументов}{\textbf{Функция, сохраняющая единицу} — функция, возвращающая единицу на единичном наборе аргументов $(11\dots 1)$. Число функций:
\BigDisp{2^{2^{n}-1}.}
\textbf{Функция, сохраняющая ноль} — функция, возвращающая ноль на нулевом наборе аргументов $(00\dots 0)$. Число функций:
\BigDisp{2^{2^{n}-1}.}}

\DefItem{Формулировка критерия Поста}{Набор (система) функций $A$ полон тогда и только тогда, когда он не содержится целиком ни в одном из предполных классов. То есть в нём есть хотя бы одна несамодвойственная, хотя бы одна нелинейная, хотя бы одна немонотонная, хотя бы одна не сохраняющая ноль и хотя бы одна не сохраняющая единицу функция.}

\DefItem{Перечислить полные системы из одной функции}{Стрелка Пирса, штрих Шеффера.}

\DefItem{Частично определенная (недоопределенная) функция}{Булева функция, значения которой определено не на всех наборах.}

\DefItem{Кодирование}{Перевод информации в удобную для передачи, обработки или хранения форму с помощью некоторого кода.}

\DefItem{Кодирование (как функция)}{Функция $c: A^* \to B^*$, где $A$ — кодируемый алфавит, $B$ — кодирующий алфавит.}

\DefItem{Однозначно декодируемый код}{Схема алфавитного кодирования называется однозначно декодируемой, если любое закодированное слово допускает только один способ декодирования.}

\DefItem{Однозначно декодируемый код (критерий)}{Функция $c$, которая обладает свойством: если
\BigDisp{c(a_1)c(a_2)\cdots c(a_k)=c(b_1)c(b_2)\cdots c(b_m),}
то $k=m$ и $a_i=b_i$ для всех $i$.}

\DefItem{Равномерный код}{Кодирование каждого символа алфавита производится с помощью строк одинаковой длины.}

\DefItem{Равномерная схема алфавитного кодирования}{Все её кодовые слова попарно различны, но имеют одинаковую длину.}

\DefItem{Энтропия сообщения}{\textbf{Формула Шеннона (информационная двоичная энтропия)} — количество информации, которое в среднем приходится на одно событие (символ):
\BigDisp{H(p_1,p_2,\dots,p_n)= -\sum_{i=1}^{n} p_i \log_2 p_i,}
где $p_i$ — вероятность $i$-го события. Эта величина также называется средней энтропией сообщения и означает измеряемое в битах среднее количество информации на символ передаваемого сообщения.}

\DefItem{Неравенство Крафта—Макмиллана}{Пусть заданы кодируемый и кодирующий алфавиты, состоящие из $n$ и $d$ символов соответственно, и заданы желаемые длины кодовых слов: $l_1,l_2,\dots,l_n$. Тогда необходимым и достаточным условием существования разделимого и префиксного кодов, обладающих заданным набором длин кодовых слов, является выполнение неравенства:
\BigDisp{\sum_{i=1}^{n} d^{-l_i}\le 1.}
Суть: префиксность не даёт выбирать слишком короткие кодовые слова; неравенство ограничивает снизу длины кодовых слов префиксного кода.}

\clearpage

\DefItem{Префиксная / суффиксная схема}{Код, в котором никакое кодовое слово не является началом другого (условие Фано). Аналогично определяется постфиксный (суффиксный) код — это код, в котором никакое кодовое слово не является концом другого. Префиксный код будет однозначно декодируемым кодом.}

\DefItem{Избыточность схемы алфавитного кодирования}{Средняя длина её кодового слова.}

\DefItem{Формула подсчета избыточности схемы}{Средняя длина кодового слова:
\BigDisp{\bar l=\sum_{i=1}^{n} p_i l_i.}}

\DefItem{Код Хаффмана}{Классический алгоритм Хаффмана на входе получает таблицу (список) частотностей символов в сообщении; далее на основании этой таблицы строится дерево кодирования Хаффмана.}

\DefItem{Код Шеннона--Фано}{Символы первичного алфавита выписывают по убыванию вероятностей; затем делят на две части с максимально близкими суммарными вероятностями.}

\DefItem{Арифметическое кодирование}{Алгоритм сжатия информации без потерь, который при кодировании ставит в соответствие тексту вещественное число из отрезка $[0;1)$.}

\DefItem{Расстояние Хэмминга}{Число позиций, в которых различаются соответствующие символы двух строк одинаковой длины.}

\end{enumerate}

\end{document}
